\chapter{Cardiopulmonary Resuscitation (CPR)}

\section*{What Is This Test or Procedure?}
CPR is an emergency procedure that combines chest compressions and artificial ventilation to manually preserve brain function and circulation in a person in cardiac arrest.

\section*{Alternative Names}
CPR, Basic Life Support, BLS

\section*{Principle}
Rhythmic chest compressions squeeze the heart against the spine, generating artificial blood flow to the brain and vital organs, while rescue breaths provide oxygenation. Maintaining perfusion during cardiac arrest bridges the patient until defibrillation and advanced care restore a spontaneous rhythm.
\section*{History}
Modern CPR was developed in the late 1950s. Kouwenhoven, Knickerbocker, and Jude discovered the efficacy of external chest compressions in 1960.

\section*{Why This Test or Procedure Is Ordered}
Performed immediately on any individual who is unresponsive, pulseless, and not breathing (cardiac arrest).

\section*{Is This a Primary or Secondary Test or Procedure}
Primary emergency resuscitation method for out-of-hospital and in-hospital cardiac arrest.

\section*{How This Test or Procedure Is Performed}
Consists of cycles of 30 high-quality chest compressions (100-120 per minute, at least 2 inches deep) followed by 2 rescue breaths. An Automated External Defibrillator (AED) is applied as soon as available.

\section*{What Preparation Is Needed}
Ensure scene safety. Assess responsiveness and breathing. Call for emergency help immediately.

\section*{Contraindications and Precautions}
Active Do Not Resuscitate (DNR) order, or obvious signs of death (rigor mortis, decapitation).

\section*{Risks}
Rib fractures, sternal fractures, gastric inflation, or aspiration of gastric contents.

\section*{What Happens After the Test or Procedure}
If Return of Spontaneous Circulation (ROSC) is achieved, immediately transition the patient to advanced post-cardiac arrest care (ventilator, targeted temperature management).

\section*{Understanding Your Results}
ROSC is identified by a palpable pulse, measurable blood pressure, and organized electrical activity on the cardiac monitor.

\section*{Normal Results}
Return of Spontaneous Circulation (ROSC); return of spontaneous respiratory effort; hemodynamically stable.

\section*{Follow-up Tests and Procedures}
Intensive care unit admission, coronary angiography, or therapeutic hypothermia.

\section*{References}
\begin{enumerate}
  \item MedlinePlus. Cardiopulmonary Resuscitation (CPR). U.S. National Library of Medicine; 2025.
  \item Current Medical Diagnosis and Treatment. McGraw-Hill; 2025.
\end{enumerate}

\clearpage
